problem:  
**Problem (Rigidity threshold for symmetry rank in full‑holonomy Calabi–Yau manifolds).**  
Let \((M^{2n},g,J,\Omega)\) be a complete, simply connected Calabi–Yau manifold with  
\[
\nabla J=0,\qquad \mathrm{Hol}(g)=SU(n),\qquad \mathrm{Ric}(g)=0.
\]  
Let \(\mathfrak{isom}(M,g)\) denote the Lie algebra of Killing fields of \(g\), and let \(\mathrm{srank}(g)\) be its **symmetry rank**, i.e. the maximal dimension of a torus that acts effectively by isometries on \((M,g)\).

**Open problem.**  
Does there exist a universal constant \(r_{\max}(n)\) (depending only on the complex dimension \(n\)) such that every **nonflat, full‑holonomy** Calabi–Yau \(n\)-fold satisfies  
\[
\mathrm{srank}(g)\le r_{\max}(n),
\]  
and equality forces a geometric model of cohomogeneity one?

In particular:
1. Determine whether there can exist full‑holonomy Calabi–Yau 4‑folds with continuous symmetry rank \(\ge 3\).  
2. More generally, characterize all maximal‑symmetry configurations consistent with \(\mathrm{Hol}(g)=SU(n)\).  
3. Does large symmetry rank inevitably reduce holonomy (e.g., to \(SU(k)\times SU(\ell)\) or \(Sp(k)\)) except in the flat or product cases?

Equivalently: **Quantify and classify the precise “symmetry threshold” beyond which a Calabi–Yau metric can no longer sustain full holonomy.**

---

Why is it a "good" problem:  

1. **Conceptually simple yet profound.**  
   The question—*“How symmetric can a Calabi–Yau manifold be before its holonomy reduces?”*—is strikingly direct, yet it probes deep structural constraints between isometry algebra, curvature, and parallel differential forms. The notion of a symmetry threshold parallels results in exceptional holonomy but remains almost completely unexplored for \(SU(n)\) holonomy beyond trivial cases.

2. **Bridges distinct frameworks.**  
   It links  
   - **holonomy theory** (Berger classification, parallel spinors),  
   - **Lie group actions** on Riemannian manifolds (isometric torus actions, cohomogeneity theory), and  
   - **Kähler geometry** (Ricci‑flatness, Kähler potentials, Hamiltonian Killing symmetries).  
   Understanding how these interact to constrain possible symmetry ranks would fuse analytic, topological, and representation‑theoretic perspectives.

3. **Truly open and quantitative.**  
   For Calabi–Yau metrics, essentially no results quantify isometric symmetry dimensions under the *full holonomy* condition. Known examples—such as the Stenzel and Candelas–de la Ossa metrics—exhibit small, specific symmetry groups, but no theorem provides universal upper bounds on symmetry rank. Even conjectural estimates could reveal new rigidity mechanisms analogous to those known for \(G_2\) and \(\mathrm{Spin}(7)\) holonomy.

4. **Potentially unifying principle.**  
   A resolution could produce a **holonomy–symmetry rigidity theorem** across all special holonomy types, illustrating that maximal symmetry and full exceptional holonomy are mutually exclusive except in codified limiting cases. Such results would clarify the geometric spectrum between flat, reducible, and genuinely irreducible Ricci‑flat geometries.

5. **Simplicity hiding depth.**  
   The problem is expressible in one line—*“Is there an upper bound on symmetry rank for full‑holonomy Calabi–Yau metrics?”*—yet solving it demands an intricate blend of analytic (Bochner and moment‑map arguments), topological (fixed‑point theorems), and algebraic‑geometric (holomorphic vector field) techniques.  
   Its conceptual purity, quantitative nature, and cross‑field depth embody the ideal “narrow entrance, profound depth” character of a genuinely good research problem.